\documentclass[10pt,a4paper]{amsart}
\usepackage[margin=19mm]{geometry}
\usepackage{amsmath,amssymb,mathtools}
\usepackage[scr=rsfs,cal=euler]{mathalfa}
\usepackage{xfrac}
\usepackage[unicode,hidelinks,bookmarks=false]{hyperref}
\newcommand{\ie}{i.e.\ }
\newcommand{\cf}{cf.\ }
\newcommand{\resp}{respectively\ }
\newcommand{\Subsection}{\subsection}
\providecommand{\nobreakdash}{\nobreak\hbox{-}}
\setlength{\emergencystretch}{2em}
\multlinegap0pt
\usepackage{amssymb}
\usepackage{xcolor}
\usepackage{xfrac}
\usepackage[shortlabels]{enumitem}
\usepackage{tikz}
\usepackage{bm}
\newtheorem{prop}{Proposition}
\newtheorem{lem}{Lemma}
\newcommand{\M}{\mathcal{M}}
\newcommand{\RR}{\mathbb{R}}
\newcommand{\cent}{\mathrm{cent}}
\title{Quantitative Approximation Rates for Group Equivariant Learning}
\author{Jonathan W. Siegel, Snir Hordan, Hannah Lawrence, Ali Syed, Nadav Dym}
\date{Source: arXiv:2602.20370v2 (CC BY 4.0)}
\begin{document}
\maketitle

\noindent\emph{Source and licence.} Quantitative Approximation Rates for Group Equivariant Learning. Version arXiv:2602.20370v2 (CC BY 4.0), distributed by arXiv under the Creative Commons Attribution 4.0 International licence, http://creativecommons.org/licenses/by/4.0/, as recorded in the arXiv licence metadata for this exact version and shown on the abstract page https://arxiv.org/abs/2602.20370v2. Full licence notice: \texttt{licenses/arxiv-2602.20370-CC-BY-4.0.txt}.\par\smallskip

\noindent\textit{Editorial scope.} Counted unit: the article's prop prop:dim (source0.tex lines 488--496). The article's complete original proof of it is delivered with it (source0.tex lines 2223--2248, one \texttt{proof} environment of 26 lines, which the article places away from the statement). No mathematics is written, repaired or re-derived: the statement and the proof are the article's own lines, in the article's own notation, and no retained line is substituted or re-flowed.  Retained uncounted as the source-local context the proof needs: lines 2022--2023; lines 2108--2110.  Nothing was omitted from any retained span, so the package carries no omission record.  No retained line contains a citation, so the wrapper carries no bibliography and no bibliography entry.  No retained line contains a figure environment, an image-inclusion command or drawing code, so no image is delivered and the package declares no asset dependency.  Editorial formatting only, in addition: an AMS \texttt{amsart} preamble that restates the article's own declarations and macros used by the retained lines, namely the environments \texttt{prop}, \texttt{lem} on separate counters, together with the standard packages those lines need.  The retained environments print in this document as Proposition 1, Lemma 1 in order of appearance, whereas the article prints the same items with the numbering of its own full document.  The complete native source of the article as distributed in the arXiv e-print for this exact version is bundled unchanged as \texttt{sources/195183/source0.tex}.
\begin{equation}\label{eq:cent}
\cent(X)=X-\frac{1}{n}X1_n1_n^T \end{equation}

\begin{lem}\label{lem:holder}
Let $V=\RR^N$ and let $G$ be a closed subgroup of the isometry group of $V$. Let $(Y,d_Y)$ be some metric space,  let $f:V\to Y $ be an invariant function and let $\tilde f:V/G\to Y$ be the induced function on the quotient space. Then $f$ is $\alpha$-H\"older if, and only if, $\tilde f$ is $\alpha$-H\"older. Moreover, $|f|_\alpha=|\tilde f|_\alpha $.
\end{lem}

\begin{prop}\label{prop:dim}[Proof in appendix]
Let  $d,n$ be natural numbers and let $G$ be a group acting on $\RR^{d\times n}$, and assume that one of the  following holds:
\begin{enumerate}
\item $G$ is the permutation group, $G=S_n$.
\item  $G=E(d)\times S_n $, and $d+1\leq n$.
\end{enumerate}
Then for any $G$ invariant set $K\subseteq \RR^{d\times n}$ for which $K/G$ is compact, 
$$N_r(K/G)\leq Cr^{-(nd-\dim(G))}, \quad \forall r>0.$$
\end{prop}

\begin{proof}[Proof of Proposition \ref{prop:dim}]
We first consider the case of \textbf{permutations}. Let $K\subseteq \RR^{d\times n}$ such that $K/G\subseteq \RR^{d\times n}/S_n $ is a compact set. The function
$$F([X])=\|X\|_F $$
is Lipschitz on $\RR^{d\times n}$ and $S_n$ invariant, and therefore it is also Lipschitz
 on the quotient space by Lemma \ref{lem:holder}. Therefore, it attains a maximum $M$ on $K/G $. Accordingly, $K/G$ is contained in $q(B_M) $, where 
$$B_M=\{X\in \RR^{d\times n}| \quad \|X\|_F\leq M \} .$$
For an appropriate $C=C(M)$, we have
$$N_r(B_M)\leq C_Mr^{-nd}, \quad \forall r>0 ,$$
and since the quotient map $q$ is 1-Lipschitz, we can push any cover of $B_M$ of balls with radius $r$ to a cover of $q(B_M)$ by balls with the same radius. Thus 
$$N_r(K/G)\leq N_r(B_M/G)\leq N_r(B_M)
\leq C_Mr^{-nd}, \quad \forall r>0 . $$

Next we consider the combined action of \textbf{permutations and rigid motions}. Let $K\subseteq \RR^{d\times n} $ be a $G$ invariant set such that $K/G$ is compact. As in \eqref{eq:cent}, we denote 
\begin{equation*}
\cent(X)=X-\frac{1}{n}X1_n1_n^T. \end{equation*}
The  function  $F([X])=\|\cent(X)\|_F$ is Lipschitz on $\RR^{d\times n}$ and is $G$ invariant. Therefore it is  well defined and Lipschitz on the quotient space, by Lemma \ref{lem:holder}, and therefore it attains a maximum $M$ on $K/G$. It follows that  $K/G $ is contained in  $q(B_M)$. Next, we claim that $\RR^{d\times n}/G $ is  contained in $q(\M)$, where  
$$\mathcal{M}=\{X\in \RR^{d\times n}| X1_n=0_d, X[i,j]=0, \forall d\geq i>j \geq 1  \} .$$
This is because, given any $X\in \RR^{d\times n}$ we can (a)  centralize $X$, and then, denoting by $k$ the rank of $\cent(X)$ (b) relabel the entries of $\tilde X=\cent(X)$ so that $\mathrm{rank}(\tilde x_1,\ldots,\tilde x_k)=k $. We can then apply the Gram-Schmidt procedure to obtain a rotation $U$ such that $U\tilde X$ is in $\M$ (if $k<d$ this rotation still exists, but is not unique), and $U\tilde X\in[X] $. It follows that $K/G\subseteq q(B_M\cap \M) $
Under our assumption that $n\geq d+1 $, the dimension of the subspace $\M$ is 
$$\dim(\mathcal{M})=nd-\frac{d(d-1)}{2}-d=nd-\dim(E(d))=nd-\dim(G),$$
and therefore there exists some $C>0$ such that 
$$N_r(B_M\cap \M)\leq Cr^{-(nd-\dim(G))}, \quad \forall r>0. $$
Again using the fact that $q$ is $1$-Lipschitz, we obtain 
$$N_r(K/G)\leq N_r(q(B_M\cap \M))\leq N_r(B_M\cap \M) \leq Cr^{-(nd-\dim(G))}, \quad \forall r>0. $$

 \end{proof}

\end{document}
